\section{Empirical strategy}\label{sec:strategy}

\subsection{Potential outcomes and the estimand}

Let $i$ index individuals, $g\in\{\Low,\mathrm{High}\}$ their skill group, and $t$ the quarter. The reform becomes effective in the second quarter of 2022, which I denote $t^{\ast}$;\footnote{Because the reform took effect on 1 May, the second quarter of 2022 contains one pre-reform and two post-reform months, and reported incomes refer to the month before the interview, so even May interviews report pre-reform (April) earnings. I therefore exclude 2022Q2 from the baseline as a partially treated transition quarter and restore it in a robustness check; the first post-reform quarter in the baseline is 2022Q3.} define the post-reform indicator $\Post_t=\mathbf{1}\{t\ge t^{\ast}\}$ and the treatment indicator $\Low_i=\mathbf{1}\{g_i=\Low\}$. Following the potential-outcomes framework, let $Y_{it}(1)$ and $Y_{it}(0)$ denote the outcome of individual $i$ in quarter $t$ with and without the reform being in force. The parameter of interest is the average effect of the reform on the low-skilled, relative to the counterfactual in which they had experienced the same evolution as the high-skilled,
\begin{equation}
\mathrm{ATT} \;=\; \mathbb{E}\!\left[\,Y_{it}(1)-Y_{it}(0)\;\middle|\;\Low_i=1,\;t\ge t^{\ast}\right].
\end{equation}
The data are repeated cross-sections rather than a panel, so the object is identified from comparisons of group means over time rather than from within-person changes.

Because Peru's minimum wage is national and changes on a single date, there is no variation in treatment \emph{timing}; every worker is exposed to the same policy at the same moment. Identification instead rests on variation in the \emph{intensity} of exposure across skill groups. The design is therefore a canonical two-group difference-in-differences with many periods, not a staggered-adoption design. This distinction matters for the choice of estimator. The recent econometric literature has shown that the two-way fixed-effects estimator can be badly biased when treatment timing is staggered and effects are heterogeneous \citep{goodmanbacon_2021, dechaisemartin_2020, sun_abraham_2021, callaway_santanna_2021}. With a common treatment date, however, these problems do not arise: the two-way fixed-effects estimator and its event-study analog recover a convex, properly weighted average of treatment effects \citep{roth_2023_trending, baker_larcker_wang_2022}. I nonetheless report the doubly robust estimator of \citet{santanna_zhao_2020} as a complementary building block.

\subsection{Estimating equations}

My baseline specification is the two-way fixed-effects difference-in-differences model
\begin{equation}
Y_{idt} \;=\; \beta\,(\Low_i\times\Post_t) \;+\; \gamma\,\Low_i \;+\; \lambda_t \;+\; \mu_d
\;+\; X_{it}'\theta \;+\; \varepsilon_{idt},
\label{eq:twfe}
\end{equation}
where $d$ denotes the department, $\lambda_t$ are quarter fixed effects, $\mu_d$ are department fixed effects, and $X_{it}$ is a vector of controls comprising age, its square, sex, an urban indicator, marital status, and years of education (which captures variation in attainment within skill groups). The coefficient $\beta$ is the difference-in-differences estimate. The quarter fixed effects absorb any shock common to both skill groups in a given quarter, including aggregate inflation, so that estimates in logs are invariant to whether wages are expressed in nominal or in nationally deflated real terms. Observations are weighted by the survey weights, and standard errors are clustered by department.

To trace the dynamics of the effect and to test the identifying assumption, I estimate the dynamic event-study specification
\begin{equation}
Y_{idt} \;=\; \sum_{k\neq -1}\beta_k\,\bigl(\Low_i\times\mathbf{1}\{t-t^{\ast}=k\}\bigr)
\;+\; \gamma\,\Low_i \;+\; \lambda_t \;+\; \mu_d \;+\; X_{it}'\theta \;+\; \varepsilon_{idt},
\label{eq:event}
\end{equation}
where $k$ indexes quarters relative to the reform and the quarter immediately before the reform ($k=-1$, the first quarter of 2022) is the omitted reference. The coefficients $\beta_k$ for $k<0$ measure differential pre-trends between the skill groups and provide a direct test of parallel trends; the coefficients for $k\ge 0$ trace the path of the effect.

\subsection{Identifying assumptions and threats}

Identification of the ATT from equation~\eqref{eq:twfe} requires that, absent the reform, the outcomes of low- and high-skilled workers would have followed parallel paths, conditional on the controls and fixed effects. Three threats deserve attention.

First, \emph{differential trends}. Education is correlated with sector, region, and formality, and the two skill groups could have been on diverging paths for reasons unrelated to the minimum wage. I assess this directly through the pre-reform event-study coefficients and a joint test of their equality to zero, and I probe robustness to department-specific linear trends. Because my window opens in the depth of the pandemic recovery, when the employment of low- and high-skilled workers rebounded at different rates, I pay particular attention to the first quarter of 2021 and report specifications that drop it.

Second, \emph{sensitivity to residual violations}. Even when pre-trends are statistically flat, they need not be exactly zero. I therefore report the sensitivity analysis of \citet{rambachan_roth_2023}, which asks how large a post-reform violation of parallel trends, relative to the largest violation observed before the reform, would be needed to overturn each conclusion.

Third, \emph{spillovers and general equilibrium}. If the reform displaces low-skilled workers into the informal sector, it may indirectly depress wages there, and if high-skilled workers are imperfect substitutes their outcomes could also move, contaminating the comparison group. I cannot rule this out, and I interpret the estimates as relative effects. The comparison group is not entirely unexposed: 40 percent of high-skilled private wage earners were paid below the new floor before the reform, against 57 percent of the low-skilled. To first order this imperfect exposure attenuates the estimates toward zero (the design identifies the effect of the seventeen-point exposure \emph{gap}, so it understates the effect on the fully exposed), but I acknowledge that it also makes the estimates relative objects rather than absolute ones. For this reason I complement the education contrast with a continuous regional-exposure design that does not rely on the comparison group at all (Section~\ref{sec:robustness}).

\subsection{Inference}

The treatment is assigned at a coarse level, the skill group, which raises the standard concern that conventional standard errors overstate precision \citep{bertrand_2004}. I address this in several complementary ways. My baseline clusters by department,\footnote{The 25 clusters are Peru's 24 departments plus the Constitutional Province of Callao. With few clusters, conventional cluster-robust standard errors can understate uncertainty, which motivates the additional corrections below.} which allows arbitrary within-department correlation across individuals and quarters; I report an alternative that clusters by department and skill group; and I apply the bias-reduced CR2 small-sample correction of \citet{pustejovsky_tipton_2018} to department-level collapsed cells, where it is computationally feasible. Because these corrections may still be optimistic with a two-group design, I complement them with randomization inference at the department level, permuting the regional bite of the minimum wage across departments to construct an exact null distribution for the exposure-intensity design.
